% Avsnitt 3.10, ur lectures/L03/README.md.

\section{Sammanfattning}
\label{c3:sec:review}

\needspace{6\baselineskip}
\subsection*{Det här ska du kunna nu}
\begin{itemize}
  \item Säga varför sekvensnät finns, förklara den funktionella skillnaden mellan ett D-lås och en
    D-vippa, och varför synkrona designer nästan uteslutande byggs av vippan.
  \item Läsa och skriva den synkrona processmallen, och känna igen ett register som flera vippor
    som delar klocka, reset och enable.
  \item Säga när en signaltilldelning får effekt: \code{<=} schemalägger snarare än tillämpar, så
    två tilldelningar i en klockad process bygger en kedja, och deras inbördes ordning spelar ingen
    roll.
  \item Bygga en flankdetektor av en vippa och en grind, och säga exakt vilken signal den
    producerar och hur länge.
\end{itemize}

\needspace{6\baselineskip}
\subsection*{Frågor att testa dig själv med}
\begin{enumerate}
  \item Vad är den funktionella skillnaden mellan ett D-lås och en D-vippa, och under vilket
    villkor ändrar var och en sin utgång?
  \item Varför uppdaterar synkrona kretsar sitt tillstånd vid en enda klockflank i stället för att
    kontinuerligt reagera på sina ingångar?
  \item Givet en signal och en vippa som håller dess värde från en cykel tidigare, hur kombinerar
    du dem till en puls som är hög i exakt en klockcykel när signalen går hög?
  \item Varför måste \code{clock} vara den enda signalen i en vippprocess känslighetslista, om inte
    designen har en asynkron reset?
  \item En klockad process innehåller \code{b <= a;} och därefter \code{c <= b;}. Vad håller
    \code{c} efter en klockflank, och varför är det inte \code{a}? Vad ändras om du byter plats på
    raderna, och varför?
\end{enumerate}

\needspace{6\baselineskip}
\subsection*{Blicken framåt}
\Chapref{c4:ch} går vidare med:
\begin{itemize}
  \item Varför den råa knappen som matar det här kapitlets flankdetektor kan driva en vippa in i
    ett odefinierat, instabilt tillstånd, och varför det är ett annat sorts problem än något här.
  \item Dubbelvippsynkroniseraren, och vad den gör och inte gör åt en studsande knapp.
  \item Samma lysdiodsväxling, ombyggd så att den faktiskt är säker på riktig hårdvara.
\end{itemize}
